\chapter{Álgebra lineal}\label{ch:b2:linalg}

El álgebra lineal del volumen del primer año trabajaba sobre $\R$ o
$\C$ en dimensión finita y admitía el \omterm{def:b2:linalg:det}{determinante} general. Este
capítulo levanta las tres restricciones: la teoría se enuncia sobre un
cuerpo $K$ arbitrario, se desarrolla sistemáticamente el juego entre un
espacio y su \emph{\omterm{def:b2:linalg:dual}{dual}} (\omterm{def:b2:linalg:dual}{bases duales}, \omterm{def:b2:linalg:annihilator}{anuladores},
\omterm{def:b2:linalg:transpose}{traspuestas}) y el \omterm{def:b2:linalg:det}{determinante} queda por fin \emph{construido} a
partir de las formas multilineales \omterm{def:b2:linalg:alternating}{alternadas} y de la signatura del~\cref{ch:b2:structures}, saldando así todas las deudas del primer año.

En todo el capítulo, $K$ es un cuerpo ($\Q$, $\R$, $\C$ o $\Z/p\Z$; a
la teoría le da igual) y, salvo mención expresa, los espacios son de
dimensión finita sobre $K$. Los resultados del primer año (bases,
dimensión, teorema del rango, matrices) se trasladan literalmente: sus
demostraciones nunca usaron nada más que los axiomas de cuerpo.

\section{Espacio dual}

\begin{definition}[Espacio dual, base dual]\label{def:b2:linalg:dual}
El \emph{dual}\index{espacio dual} de $E$ es $E^* = \mathcal{L}(E,
K)$, el espacio de las formas lineales. Si $\mathcal{B} = (e_1, \dots,
e_n)$ es una base de $E$, las \emph{formas coordenadas} $e_1^*, \dots,
e_n^*$ definidas por $e_i^*(e_j) = \delta_{ij}$ (delta de Kronecker:
$1$ si $i = j$, y $0$ en otro caso) forman la \emph{base
dual}\index{base dual} $\mathcal{B}^*$ de $E^*$; en particular $\dim
E^* = \dim E$, y
\[
x = \sum_{i=1}^{n} e_i^*(x)\, e_i \quad (x \in E),
\qquad
\varphi = \sum_{i=1}^{n} \varphi(e_i)\, e_i^* \quad (\varphi \in
E^*).
\]
\end{definition}

\begin{proof}[Demostración de que $\mathcal{B}^*$ es una base]
Libre: aplicando una combinación nula $\sum \lambda_i e_i^* = 0$ a
$e_j$ se obtiene $\lambda_j = 0$. Generadora: para $\varphi \in E^*$,
la forma $\varphi - \sum_i \varphi(e_i) e_i^*$ anula todos los $e_j$ y
por tanto es nula (una aplicación lineal que se anula sobre una base
es nula). Las dos fórmulas destacadas son esos mismos cálculos leídos
en sentido directo.
\end{proof}

\begin{example}\label{ex:b2:linalg:dualexamples}
En $K_n[X]$ con la base $(1, X, \dots, X^n)$, la \omterm{def:b2:linalg:dual}{base dual} es $P
\mapsto \frac{P^{(k)}(0)}{k!}$ (coeficientes de Taylor). Otra base del
\omterm{def:b2:linalg:dual}{dual}: las evaluaciones $P \mapsto P(x_i)$ en $n + 1$ puntos
distintos; su base ``predual'' en $K_n[X]$ es exactamente la familia
de los polinomios de Lagrange $L_i$ (volumen del primer año), ya que
$L_i(x_j) = \delta_{ij}$. La interpolación \emph{es} dualidad.
\end{example}

\begin{method}[Bases duales y anteduales en la práctica]\label{met:b2:linalg:dualbasis}
Para desarrollar una forma $\varphi$ sobre una base $(e_i)$ de $E$:
las coordenadas son los \emph{valores} $\varphi(e_i)$; no hay ningún
sistema que resolver. Para hallar la base $(u_j)$ de $E$ cuya
\omterm{def:b2:linalg:dual}{dual} es una base dada $(\varphi_1, \dots, \varphi_n)$ de $E^*$
(la \emph{antedual}), se resuelven los $n$ sistemas lineales
\[
\varphi_i(u_j) = \delta_{ij} \qquad (1 \leq i \leq n),
\]
una columna $u_j$ cada vez; en términos matriciales, si las filas
de $M$ dan los coeficientes de los $\varphi_i$ en una base conocida
de $E^*$, las columnas de $M^{-1}$ son los $u_j$. La existencia y
la unicidad de la antedual se demuestran en el problema de fin de
semana de este capítulo; el cálculo es siempre esta inversión.
\end{method}

\begin{example}[Una base dual de $\R^2$, calculada por completo]\label{ex:b2:linalg:dualbasis2d}
Para la base $b_1 = (1, 1)$, $b_2 = (1, -1)$ de $\R^2$, la
\omterm{def:b2:linalg:dual}{base dual} $(b_1^*, b_2^*)$ ha de cumplir $b_i^*(b_j) =
\delta_{ij}$. Escribiendo $b_1^*(x, y) = \alpha x + \beta y$, las
condiciones $\alpha + \beta = 1$ y $\alpha - \beta = 0$ dan
\[
b_1^*(x, y) = \frac{x + y}{2},
\qquad\text{y análogamente}\qquad
b_2^*(x, y) = \frac{x - y}{2} .
\]
Comprobaciones de sensatez: $b_1^*$ \emph{no} es $e_1^* + e_2^*$
evaluado ingenuamente; la \omterm{def:b2:linalg:dual}{base dual} depende de la base entera y no
de cada vector por separado (al sustituir $b_2$ por $(0, 1)$,
$b_1^*$ pasa a ser $x \mapsto x$). Y la fórmula de desarrollo
funciona: $(x, y) = \frac{x+y}2\,b_1 + \frac{x-y}2\,b_2$, la
descomposición en parte par e impar de un par; las \omterm{def:b2:linalg:dual}{bases duales}
son extractores de coordenadas, y esta extrae las partes simétrica
y antisimétrica.
\end{example}

\begin{definition}[Anulador]\label{def:b2:linalg:annihilator}
Para un subespacio $F \subseteq E$, el
\emph{anulador}\index{anulador} es
\[
F^{\circ} = \{\varphi \in E^* : \varphi|_F = 0\},
\]
un subespacio de $E^*$.
\end{definition}

\begin{theorem}[Dimensión del anulador]\label{thm:b2:linalg:annihilator}
$\dim F^{\circ} = \dim E - \dim F$. Además, $F \mapsto F^\circ$
invierte las inclusiones, y $F$ se recupera a partir de su \omterm{def:b2:linalg:annihilator}{anulador}:
\[
F = \{x \in E : \forall\varphi \in F^\circ,\ \varphi(x) = 0\}.
\]
En consecuencia, todo subespacio de dimensión $p$ en dimensión $n$ es
el conjunto de soluciones de $n - p$ ecuaciones lineales
independientes, y recíprocamente.
\end{theorem}

\begin{proof}
Elijamos una base $(e_1, \dots, e_p)$ de $F$ completada hasta una base
de $E$. Una forma $\varphi = \sum \varphi(e_i) e_i^*$ anula $F$ si y
solo si se anulan sus $p$ primeros coeficientes: $F^\circ =
\operatorname{Vect}(e_{p+1}^*, \dots, e_n^*)$, de dimensión $n - p$.
La inversión de las inclusiones es inmediata. Para la recuperación:
el miembro derecho contiene a $F$; recíprocamente, si $x \notin F$,
complétese una base de $F$ con $x$ y con otros vectores; la forma
coordenada de $x$ en esa base anula $F$ pero no $x$. La lectura en
términos de ``ecuaciones'' se obtiene tomando una base $(\varphi_1,
\dots, \varphi_{n-p})$ de $F^\circ$: entonces $F = \bigcap
\ker\varphi_j$, intersección de $n - p$ hiperplanos independientes.
\end{proof}

\begin{example}[Un anulador, en los dos sentidos]\label{ex:b2:linalg:annihilatorwork}
Sea $F = \operatorname{Vect}\bigl((1, 2, 1),\ (1, 0, -1)\bigr)
\subseteq \R^3$. Una forma $\varphi = a\,e_1^* + b\,e_2^* +
c\,e_3^*$ anula $F$ si y solo si
\[
a + 2b + c = 0
\qquad\text{y}\qquad
a - c = 0 ,
\]
es decir, $c = a$ y $b = -a$: $F^\circ = \R\,(e_1^* - e_2^* +
e_3^*)$, de dimensión $3 - 2 = 1$, como exige el~\cref{thm:b2:linalg:annihilator}. Leído al revés: $F = \{(x, y, z)
: x - y + z = 0\}$, el plano recuperado como núcleo de la única
forma que genera $F^\circ$. Pasar de una familia generadora a unas
ecuaciones \emph{es} calcular un \omterm{def:b2:linalg:annihilator}{anulador}; pasar de las ecuaciones a
una parametrización es calcular un preanulador. (Comprobación:
ambos vectores generadores cumplen $x - y + z = 0$.)
\end{example}

\begin{definition}[Aplicación traspuesta]\label{def:b2:linalg:transpose}
Para $u \in \mathcal{L}(E, F)$, la \emph{traspuesta}\index{aplicación
traspuesta} $u^{\mathsf T} \in \mathcal{L}(F^*, E^*)$ es
\[
u^{\mathsf T}(\psi) = \psi \circ u .
\]
Cumple $(v \circ u)^{\mathsf T} = u^{\mathsf T} \circ v^{\mathsf T}$ y,
en \omterm{def:b2:linalg:dual}{bases duales}, la matriz de $u^{\mathsf T}$ es la matriz
traspuesta de la de $u$, lo que por fin \emph{explica} la traspuesta
del primer año.
\end{definition}

\begin{example}[La traspuesta, entrada a entrada]\label{ex:b2:linalg:transposework}
Sea $u \colon \R^2 \to \R^3$ de matriz $A =
\begin{psmallmatrix} 1 & 2\\ 0 & 1\\ 3 & 0\end{psmallmatrix}$ en
las bases canónicas. Para $\psi = b_1f_1^* + b_2f_2^* + b_3f_3^*
\in (\R^3)^*$, calculemos $u^{\mathsf T}(\psi) = \psi \circ u$
sobre la base de $\R^2$:
\[
(\psi \circ u)(e_1) = \psi(1, 0, 3) = b_1 + 3b_3,
\qquad
(\psi \circ u)(e_2) = \psi(2, 1, 0) = 2b_1 + b_2 .
\]
Luego $u^{\mathsf T}(\psi) = (b_1 + 3b_3)\,e_1^* + (2b_1 +
b_2)\,e_2^*$, y en las \omterm{def:b2:linalg:dual}{bases duales} la matriz de $u^{\mathsf T}$
es
\[
\begin{pmatrix} 1 & 0 & 3\\ 2 & 1 & 0\end{pmatrix}
= A^{\mathsf T} :
\]
la \omterm{def:b2:linalg:transpose}{traspuesta} abstracta \emph{es} la matriz volteada, sin dejar
ningún cálculo a la fe. Obsérvese el mecanismo: la \emph{columna}
$j$-ésima de $A$ pasó a ser la \emph{fila} $j$-ésima de la nueva
matriz porque $\psi \circ u$ lee las salidas de $u$ a través de los
coeficientes de $\psi$.
\end{example}

\begin{proposition}\label{prop:b2:linalg:transposerank}
$\ker u^{\mathsf T} = (\operatorname{im} u)^{\circ}$ y
$\operatorname{im} u^{\mathsf T} = (\ker u)^{\circ}$. En
consecuencia, $\operatorname{rk}(u^{\mathsf T}) =
\operatorname{rk}(u)$: el rango por filas es igual al rango por
columnas, demostrado estructuralmente.
\end{proposition}

\begin{proof}
$\psi \in \ker u^{\mathsf T} \iff \psi \circ u = 0 \iff \psi$ anula
$\operatorname{im} u$: esa es la primera identidad. Para la segunda:
$u^{\mathsf T}(\psi) = \psi \circ u$ anula siempre $\ker u$, luego
$\operatorname{im} u^{\mathsf T} \subseteq (\ker u)^\circ$; las
dimensiones coinciden por el teorema del rango y el~\cref{thm:b2:linalg:annihilator}:
\[
\operatorname{rk} u^{\mathsf T} = \dim F^* - \dim\ker u^{\mathsf T}
= \dim F - \bigl(\dim F - \operatorname{rk} u\bigr)
= \operatorname{rk} u
= \dim (\ker u)^{\circ} . \qedhere
\]
\end{proof}

\begin{example}[El rango leído por los dos lados]\label{ex:b2:linalg:rankbothsides}
Sea
\[
A = \begin{pmatrix}
1 & 2 & 0 & 1\\
0 & 1 & 1 & 1\\
1 & 3 & 1 & 2
\end{pmatrix} .
\]
\emph{Rango por columnas:} la tercera fila es la suma de las dos
primeras, luego $\operatorname{rk} A \leq 2$; las columnas $1$ y $2$
son libres: $\operatorname{rk} A = 2$. \emph{Núcleo de la
\omterm{def:b2:linalg:transpose}{traspuesta}:} resolviendo $A^{\mathsf T}y = 0$ se obtiene $y \in
\R\,(1, 1, -1)$, de modo que $\ker A^{\mathsf T}$ tiene dimensión $1
= 3 - 2$: exactamente $(\operatorname{im} A)^\circ$ bajo la
identificación de $(\R^3)^*$ con los vectores fila, tal como afirma
la~\cref{prop:b2:linalg:transposerank}; la única relación
``fila$_3$ = fila$_1$ + fila$_2$'' \emph{es} el \omterm{def:b2:linalg:annihilator}{anulador} del
espacio de columnas. El rango por filas ($2$ filas libres) y el
rango por columnas coinciden no por casualidad, sino porque ambos
valen $\operatorname{rk} A = \operatorname{rk} A^{\mathsf T}$.
\end{example}

\begin{example}[La dualidad lee una fórmula de cuadratura]\label{ex:b2:linalg:quadratureread}
¿Por qué existe una regla como la de Simpson
(\cref{exo:b2:linalg:4}) y por qué es única? La dualidad responde
antes de todo cálculo. Sobre $E = \R_2[X]$, la integral $P \mapsto
\int_0^1 P$ es un vector concreto del \omterm{def:b2:linalg:dual}{dual} $E^*$, de dimensión
tres; las evaluaciones en $0$, $\frac12$ y $1$ forman una
\emph{base} de $E^*$; por tanto la integral se desarrolla de manera
única sobre ellas, y ese desarrollo \emph{es} la regla de Simpson,
con sus coeficientes incluidos. Un recuento de dimensiones también
calibra las expectativas: sobre $\R_3[X]$, cuatro dimensiones de
formas no pueden en general ser \omterm{def:b2:structures:generated}{generadas} por tres evaluaciones, de
modo que la exactitud sobre las cúbicas no se la debemos a la
dualidad; que Simpson integre las cúbicas exactamente pese a todo
es una simetría de regalo (cancelación de los grados impares en
torno a $\frac12$), y hay que comprobarla a mano. Las reglas con $n
+ 1$ nodos son desarrollos de la forma integración en una base de
evaluaciones de $\R_n[X]^*$: la existencia y la unicidad cuestan un
teorema sobre \omterm{def:b2:linalg:dual}{bases duales}; solo los grados de regalo cuestan
trabajo.
\end{example}

\section{Formas multilineales alternadas}

\begin{definition}\label{def:b2:linalg:alternating}
Una aplicación $f \colon E^n \to K$ es \emph{$n$-lineal} cuando es
lineal en cada variable, y \emph{alternada}\index{forma alternada}
cuando se anula siempre que dos argumentos son iguales. Ser alternada
implica ser \emph{antisimétrica}: intercambiar dos argumentos cambia
el signo (desarróllese $f(\dots, x + y, \dots, x + y, \dots) = 0$);
más en general, para $\sigma \in \mathfrak{S}_n$,
\[
f(x_{\sigma(1)}, \dots, x_{\sigma(n)}) =
\varepsilon(\sigma)\, f(x_1, \dots, x_n),
\]
descomponiendo $\sigma$ en \omterm{def:b2:structures:sn}{transposiciones}
(\cref{thm:b2:structures:signature}).
\end{definition}

\begin{theorem}[Teorema fundamental de los determinantes]\label{thm:b2:linalg:detspace}
Sea $\dim E = n$ y $\mathcal{B} = (e_1, \dots, e_n)$ una base. El
espacio de las formas $n$-lineales \omterm{def:b2:linalg:alternating}{alternadas} sobre $E$ tiene
dimensión $1$: toda forma así es múltiplo de
\[
\det{}_{\mathcal{B}}(x_1, \dots, x_n)
= \sum_{\sigma \in \mathfrak{S}_n} \varepsilon(\sigma)
\prod_{i=1}^{n} a_{\sigma(i),\,i},
\qquad
x_j = \sum_{i} a_{ij} e_i ,
\]
y $\det_{\mathcal{B}}$ es la única que toma el valor $1$ sobre
$\mathcal{B}$.\index{determinante!construcción}
\end{theorem}

\begin{proof}
Sea $f$ una forma $n$-lineal \omterm{def:b2:linalg:alternating}{alternada}. Desarrollando cada argumento
sobre $\mathcal{B}$ por multilinealidad,
\[
f(x_1, \dots, x_n)
= \sum_{i_1, \dots, i_n} a_{i_1,1}\cdots a_{i_n,n}\,
f(e_{i_1}, \dots, e_{i_n}).
\]
Los términos con un índice repetido se anulan (por ser \omterm{def:b2:linalg:alternating}{alternada});
las tuplas $(i_1, \dots, i_n)$ que sobreviven son las inyectivas, es
decir, $i_k = \sigma(k)$ para una permutación $\sigma$, y la
antisimetría reordena $f(e_{\sigma(1)}, \dots, e_{\sigma(n)}) =
\varepsilon(\sigma) f(e_1, \dots, e_n)$. Por tanto
\[
f = f(e_1, \dots, e_n) \cdot \det{}_{\mathcal{B}} :
\]
toda \omterm{def:b2:linalg:alternating}{forma alternada} es ese múltiplo, siempre que el propio
$\det_{\mathcal{B}}$ (la suma destacada) \emph{sea} $n$-lineal
\omterm{def:b2:linalg:alternating}{alternada} y tome el valor $1$ sobre $\mathcal B$. La
multilinealidad es clara (cada sumando es lineal en cada columna).
Valor sobre $\mathcal B$: el único término no nulo es el de $\sigma
= \mathrm{id}$. \omterm{def:b2:linalg:alternating}{Alternada}: supongamos $x_j = x_k$ ($j \neq k$),
de modo que las columnas de coordenadas cumplen $a_{i j} = a_{i k}$
para todo $i$. Emparejemos cada $\sigma$ con $\sigma' =
\sigma\circ(j\,k)$, una involución sin puntos fijos de
$\mathfrak{S}_n$. Los productos emparejados coinciden:
\[
\prod_i a_{\sigma'(i),\,i}
= a_{\sigma(k),\,j}\; a_{\sigma(j),\,k}
\prod_{i \neq j,k} a_{\sigma(i),\,i}
= a_{\sigma(k),\,k}\; a_{\sigma(j),\,j}
\prod_{i \neq j,k} a_{\sigma(i),\,i}
= \prod_i a_{\sigma(i),\,i},
\]
usando la igualdad de las columnas $j$ y $k$; mientras que
$\varepsilon(\sigma') = -\varepsilon(\sigma)$. Cada pareja aporta
cero: la suma se anula.
\end{proof}

\begin{example}[Sarrus, deducida y demolida]\label{ex:b2:linalg:sarrus}
Para $n = 3$, la fórmula de las permutaciones tiene exactamente $3!
= 6$ términos. Listando $\mathfrak{S}_3$ por signatura
---$\mathrm{id}$, $(1\,2\,3)$, $(1\,3\,2)$ pares; $(1\,2)$,
$(1\,3)$, $(2\,3)$ impares--- se obtiene
\[
\det A = a_{11}a_{22}a_{33} + a_{21}a_{32}a_{13} +
a_{31}a_{12}a_{23}
- a_{21}a_{12}a_{33} - a_{31}a_{22}a_{13} -
a_{11}a_{32}a_{23} :
\]
precisamente la regla de las ``diagonales'' de Sarrus que se
enseña en secundaria, ahora un teorema y con los signos
misteriosos identificados como signaturas. La demolición: para $n
= 4$ hay $24$ permutaciones, de las cuales cualquier esquema de
diagonales solo recoge $8$; Sarrus no tiene versión en grado $4$, y
el desarrollo por cofactores (\cref{thm:b2:linalg:detrules}~(4))
toma el relevo. Contar términos es además una advertencia: la
fórmula de las permutaciones tiene $n!$ sumandos, así que es una
\emph{definición}, no un algoritmo; el método de eliminación
calcula $\det$ en $O(n^3)$ operaciones.
\end{example}

\begin{definition}[Determinantes]\label{def:b2:linalg:det}
El \emph{determinante de una familia}\index{determinante} en una base
es $\det_{\mathcal{B}}(x_1, \dots, x_n)$; el \emph{determinante de una
matriz} $A$ es el determinante de sus columnas en la base canónica ---
la fórmula de las permutaciones de más arriba---; el
\emph{determinante de un endomorfismo} $u$ es el escalar $\det u$ tal
que
\[
\det{}_{\mathcal{B}}\bigl(u(x_1), \dots, u(x_n)\bigr)
= \det u \cdot \det{}_{\mathcal{B}}(x_1, \dots, x_n)
\quad \text{para todos los } x_i
\]
(el miembro izquierdo es $n$-lineal alternado, luego múltiplo de
$\det_\mathcal{B}$ por el~\cref{thm:b2:linalg:detspace}; el factor no
depende de $\mathcal{B}$).
\end{definition}

\begin{theorem}[El cálculo de determinantes, demostrado]\label{thm:b2:linalg:detrules}
\begin{enumerate}
  \item $\det(uv) = \det u\,\det v$; $\;\det(AB) = \det A \det B$.
  \item $u$ es invertible $\iff \det u \neq 0$; una familia es base
        $\iff$ su \omterm{def:b2:linalg:det}{determinante} en alguna base es no nulo.
  \item $\det(A^{\mathsf T}) = \det A$.
  \item El desarrollo por cofactores a lo largo de cualquier fila o
        columna, tal como se enunció en el volumen del primer año, es
        válido; y dos matrices semejantes tienen el mismo
        \omterm{def:b2:linalg:det}{determinante}.
\end{enumerate}
\end{theorem}

\begin{proof}
(1) Aplíquese dos veces la relación de definición:
$\det_{\mathcal B}(uv(x_i)) = \det u \cdot \det_{\mathcal
B}(v(x_i)) = \det u \det v \cdot \det_{\mathcal B}(x_i)$.

(2) Si $u$ es invertible, $\det u \det u^{-1} = \det \mathrm{id} =
1 \neq 0$. Si no lo es, las imágenes $u(e_i)$ están ligadas;
expresando una a través de las demás y desarrollando,
$\det_{\mathcal B}(u(e_i)) = 0$ (el carácter alternado mata las
direcciones repetidas), luego $\det u = 0$. El criterio de base es
el mismo enunciado para familias.

(3) En la fórmula de las permutaciones, reindexemos cada producto
mediante $j = \sigma(i)$, es decir, $i = \tau(j)$ con $\tau =
\sigma^{-1}$: los factores son los mismos números en otro orden,
así que
\[
\prod_{i=1}^{n} a_{\sigma(i),\,i} = \prod_{j=1}^{n}
a_{j,\,\tau(j)} ,
\]
y $\varepsilon(\tau) = \varepsilon(\sigma)^{-1} =
\varepsilon(\sigma)$ (los valores son $\pm1$ y $\varepsilon$ es un
morfismo). Sumar sobre $\sigma$ es lo mismo que sumar sobre $\tau$
(la inversión es una biyección de $\mathfrak{S}_n$):
\[
\det A = \sum_{\tau}\varepsilon(\tau)\prod_j a_{j,\tau(j)}
= \det(A^{\mathsf T}),
\]
siendo la última suma la fórmula de las permutaciones aplicada a
las entradas \omterm{def:b2:linalg:transpose}{traspuestas} $(A^{\mathsf T})_{ij} = a_{ji}$.

(4) Fijemos la columna $j$ y separemos $x_j = \sum_i a_{ij} e_i$ por
linealidad: $\det A = \sum_i a_{ij}\, \det(\dots, e_i, \dots)$, y
llevar $e_i$ a la última posición ($n - i$ \omterm{def:b2:structures:sn}{transposiciones} de filas
y $n - j$ de columnas, vía (3)) identifica $\det(\dots, e_i, \dots)
= (-1)^{i+j}\Delta_{ij}$ con el menor: exactamente la regla de los
cofactores del primer año. Semejanza: $\det(P^{-1}AP) = \det
P^{-1}\det A \det P = \det A$ por (1).
\end{proof}

\begin{example}[Desarrollo por cofactores, ejecutado]\label{ex:b2:linalg:cofactorwork}
Calculemos
\[
\det\begin{pmatrix}
2 & 1 & 3\\
0 & 4 & 1\\
1 & 2 & 0
\end{pmatrix}
\]
a lo largo de la primera columna (con la pereza que permite un
cero: uno). Los signos siguen el tablero $(-1)^{i+j}$:
\[
2\,\det\begin{pmatrix}4 & 1\\ 2 & 0\end{pmatrix}
- 0
+ 1\cdot\det\begin{pmatrix}1 & 3\\ 4 & 1\end{pmatrix}
= 2(0 - 2) + (1 - 12) = -15 .
\]
Comprobación por Sarrus (\cref{ex:b2:linalg:sarrus}): $0 + 1 + 0 -
12 - 0 - 4 = -15$. Estrategia, no dogma: desarróllese por la línea
con más ceros y, cuando ninguna tenga, créense primero mediante
operaciones elementales; una ronda de eliminación cuesta menos que
dos capas de cofactores.
\end{example}

\begin{example}[Un determinante mediante las reglas]\label{ex:b2:linalg:onesmatrix}
Sea $J \in \mathcal{M}_n(K)$ la matriz de unos y $a \in K$;
calculemos $\det(aI_n + J)$ con las herramientas recién
demostradas. Todas las columnas de $aI_n + J$ suman igual: sumemos
todas las filas a la primera (el \omterm{def:b2:linalg:det}{determinante} no varía, pues sumar
un múltiplo de una fila a otra añade un término con una dirección
repetida, que la alternancia mata). La primera fila pasa a ser $(a
+ n, a + n, \dots, a + n)$; saquemos $a + n$ por linealidad en esa
fila y restemos después la primera columna de todas las demás: lo
que queda es triangular con diagonal $(1, a, \dots, a)$. Por tanto
\[
\det(aI_n + J) = (a + n)\,a^{\,n-1}.
\]
La moraleja: las raíces $a = 0$ (de multiplicidad $n - 1$) y $a =
-n$ dicen que $J$ tiene el valor propio $0$ con multiplicidad $n -
1$ y el valor propio $n$ una vez, es decir, el espectro de la
matriz $J$ de rango uno, un capítulo antes de tiempo (el~\cref{ch:b2:reduction} lo hará sistemático).
\end{example}

\begin{example}[Un determinante por la fórmula de las permutaciones]\label{ex:b2:linalg:permexample}
Para una matriz con muchos ceros la fórmula resulta práctica por sí
sola: en
\[
A = \begin{pmatrix}
0 & a & 0 & 0\\
0 & 0 & b & 0\\
0 & 0 & 0 & c\\
d & 0 & 0 & 0
\end{pmatrix},
\]
la única permutación que selecciona entradas no nulas es el \omterm{def:b2:structures:sn}{ciclo}
de longitud $4$ $\sigma = (1\,2\,3\,4)$, que envía la columna $1$ a
la fila $4$, etc.; como $\varepsilon(\sigma) = (-1)^3 = -1$, resulta
$\det A = -abcd$. (Compruébese mediante tres intercambios de
columnas hasta llegar a una matriz diagonal.)
\end{example}

\begin{example}[Un Vandermonde por la fórmula del producto]\label{ex:b2:linalg:vandermondework}
Para los nodos $0, 1, 2$ (usados por reglas de cuadratura como la
del~\cref{exo:b2:linalg:4}), el \omterm{def:b2:linalg:det}{determinante} de Vandermonde del~\cref{exo:b2:linalg:11} se evalúa de un vistazo:
\[
\det\begin{pmatrix}
1 & 1 & 1\\
0 & 1 & 2\\
0 & 1 & 4
\end{pmatrix}
= (1 - 0)(2 - 0)(2 - 1) = 2 ,
\]
y por desarrollo directo a lo largo de la primera columna:
$1\cdot(4 - 2) = 2$; coinciden. Que no se anule cuando los nodos
son distintos es toda la teoría de la interpolación en un solo
\omterm{def:b2:linalg:det}{determinante}: las formas de evaluación $P \mapsto P(a_i)$ son una
base del \omterm{def:b2:linalg:dual}{dual} exactamente cuando este \omterm{def:b2:linalg:det}{determinante} no es nulo, es
decir, siempre que los $a_i$ sean distintos; el~\cref{ex:b2:linalg:dualexamples}, cuantificado.
\end{example}

\section{La traza, revisitada}

\begin{proposition}\label{prop:b2:linalg:trace}
La traza $\operatorname{tr} \colon \mathcal{M}_n(K) \to K$ es la
única forma lineal con $\operatorname{tr}(AB) =
\operatorname{tr}(BA)$ y $\operatorname{tr}(I_n) = n$ (si
$\operatorname{char} K = 0$); la traza de un endomorfismo está bien
definida a través de cualquier representación matricial, y
\[
\operatorname{tr}(u) = \sum_{i} e_i^*\bigl(u(e_i)\bigr)
\]
en cualquier base: la dualidad escribe la traza sin recurrir a una
base.
\end{proposition}

\begin{proof}
Las identidades $\operatorname{tr}(AB) = \operatorname{tr}(BA)$ y la
invariancia por cambio de base se demostraron en el primer año.
Unicidad: una forma lineal $t$ con $t(AB) = t(BA)$ anula todo
conmutador $AB - BA$. Afirmamos que los conmutadores generan el
hiperplano de traza nula, de dimensión $n^2 - 1$. Bastan dos
familias de conmutadores. La regla de multiplicación de las matrices
elementales es $E_{ab}E_{cd} = \delta_{bc}E_{ad}$. Para $i \neq j$ da
\[
E_{ii}E_{ij} - E_{ij}E_{ii} = E_{ij} - 0 = E_{ij}
\]
(el segundo producto es $E_{ij}E_{ii} = \delta_{ji}E_{ii} = 0$,
pues $j \neq i$): toda $E_{ij}$ fuera de la diagonal es un
conmutador. Y
\[
E_{ij}E_{ji} - E_{ji}E_{ij} = E_{ii} - E_{jj} .
\]
Las $E_{ij}$ ($i \neq j$; hay $n^2 - n$) junto con las $E_{11} -
E_{jj}$ ($j \geq 2$; hay $n - 1$) son $n^2 - 1$ matrices de traza
nula linealmente independientes: generan el hiperplano
$\ker\operatorname{tr}$. Así pues, $t$ se anula donde lo hace
$\operatorname{tr}$ y factoriza a través de ella: $t =
c\operatorname{tr}$; después, $t(I) = n$ obliga a $c = 1$. En cuanto
a la fórmula destacada: la $i$-ésima entrada diagonal de la matriz
de $u$ es precisamente $e_i^*(u(e_i))$.
\end{proof}

\begin{remark}[Errores frecuentes]\label{rem:b2:linalg:pitfalls}
(i) El \omterm{def:b2:linalg:det}{determinante} es $n$-lineal en las \emph{columnas}, no lineal
en la matriz: $\det(A + B) \neq \det A + \det B$ en general, y
$\det(\lambda A) = \lambda^n\det A$, no $\lambda\det A$. (ii) La
trasposición invierte los productos: $(vu)^{\mathsf T} =
u^{\mathsf T}v^{\mathsf T}$; olvidar la inversión arruina todo
cálculo con inversas. (iii) El \omterm{def:b2:linalg:annihilator}{anulador} $F^\circ$ vive en $E^*$,
no en $E$: solo se convierte en el familiar ``complemento
ortogonal'' una vez que un producto escalar identifica $E$ con $E^*$
(\cref{ch:b2:quadratic}); ninguna identificación de ese tipo es
canónica. (iv) ``El rango por filas es igual al rango por columnas''
no significa que el \emph{espacio} de filas sea igual al de
columnas: ambos viven en espacios distintos ($K^n$ y $K^m$) y se
relacionan mediante la~\cref{prop:b2:linalg:transposerank}, no son
iguales. (v) La fórmula de las permutaciones es un instrumento de
demostración: para calcular con números, úsense operaciones
elementales y cofactores (\cref{ex:b2:linalg:sarrus}).
\end{remark}

\begin{example}[El emparejamiento traza parte el espacio de matrices]\label{ex:b2:linalg:tracepairing}
Sobre $\mathcal{M}_2(\R)$, con el emparejamiento $\langle A,
B\rangle = \operatorname{tr}(AB)$ del~\cref{exo:b2:linalg:9},
descompongamos $M = \begin{psmallmatrix}1 & 4\\ 2 &
3\end{psmallmatrix}$ en sus partes simétrica y antisimétrica:
\[
M = S + A, \qquad
S = \tfrac12(M + M^{\mathsf T}) =
\begin{pmatrix}1 & 3\\ 3 & 3\end{pmatrix},
\qquad
A = \tfrac12(M - M^{\mathsf T}) =
\begin{pmatrix}0 & 1\\ -1 & 0\end{pmatrix}.
\]
Entonces $\operatorname{tr}(SA) = \operatorname{tr}
\begin{psmallmatrix}-3 & 1\\ -3 & 3\end{psmallmatrix} = 0$: las dos
partes son ``ortogonales'' para el emparejamiento traza, caso
particular del hecho general (demostrado en el problema de fin de
semana de este capítulo) de que las matrices antisimétricas forman
exactamente el \omterm{def:b2:linalg:annihilator}{anulador} de las simétricas. La dualidad ve la
descomposición $\mathcal{M}_n = \mathcal{S}_n \oplus \mathcal{A}_n$
antes de que se elija ningún producto escalar.
\end{example}

\begin{remark}[Perspectivas dentro de este volumen]\label{rem:b2:linalg:perspectives}
Conviene observar cómo las tres construcciones de este capítulo se
cambian de traje más adelante. La \emph{\omterm{def:b2:linalg:transpose}{traspuesta}} vuelve en el~\cref{ch:b2:reduction}: $u$ y $u^{\mathsf T}$ comparten valores
propios con multiplicidades geométricas iguales (problema de fin de
semana de este capítulo, pregunta 15), y por eso los análisis por
filas y por columnas de una matriz nunca discrepan. El
\emph{\omterm{def:b2:linalg:det}{determinante}} pasa a ser función de un parámetro en el~\cref{ch:b2:reduction} ($\chi_u(X) = \det(X\,\mathrm{id} - u)$) y un
jacobiano en el~\cref{ch:b2:multint}, donde su multilinealidad se
convierte en el factor del cambio de variable. La \emph{traza}
siembra los invariantes de semejanza: es el segundo coeficiente de
$\chi_u$, la suma de los valores propios y, con el tiempo, la
integral de la diagonal en identidades del estilo del~\cref{ch:b2:fourier}. Un capítulo de álgebra lineal, tres largas
sombras.
\end{remark}

\begin{remark}[Dónde se usa este capítulo]
El \omterm{def:b2:linalg:dual}{espacio dual} no es una abstracción gratuita: los \omterm{def:b2:linalg:annihilator}{anuladores} y
las \omterm{def:b2:linalg:transpose}{traspuestas} gobiernan la teoría de resolubilidad de los
sistemas lineales (el problema de fin de semana de este capítulo
demuestra a partir de ellos la alternativa de Fredholm en dimensión
finita), los emparejamientos no degenerados reaparecen como la forma
polar en el~\cref{ch:b2:quadratic} y el adjunto en el~\cref{ch:b2:hermitian}, y el \omterm{def:b2:linalg:det}{determinante} construido aquí mueve todo
el~\cref{ch:b2:reduction}. En el volumen del tercer año, esa misma
dualidad, transportada a dimensión infinita, se convierte en el
teorema de representación de Riesz y en la teoría de Fredholm sobre
espacios de Hilbert, con la compacidad ocupando el lugar de los
recuentos de dimensiones usados aquí.
\end{remark}

\section{Ejercicios}

\begin{exercise}[$\star$]\label{exo:b2:linalg:1}
En $\R^3$, sean $\varphi_1(x,y,z) = x + y$, $\varphi_2 = y + z$,
$\varphi_3 = x + z$. Demuestra que $(\varphi_1, \varphi_2,
\varphi_3)$ es una base de $(\R^3)^*$ y halla la base de $\R^3$ de la
que es la \omterm{def:b2:linalg:dual}{dual}.
\end{exercise}

\begin{exercise}[$\star$]\label{exo:b2:linalg:2}
Calcula mediante la fórmula de las permutaciones los
\omterm{def:b2:linalg:det}{determinantes} de
\[
\begin{pmatrix} 0 & 0 & a\\ 0 & b & 0\\ c & 0 & 0 \end{pmatrix},
\qquad
\begin{pmatrix}
a & b & 0 & 0\\
c & d & 0 & 0\\
0 & 0 & e & f\\
0 & 0 & g & h
\end{pmatrix},
\]
y enuncia la regla para matrices diagonales por bloques que sugiere
el segundo.
\end{exercise}

\begin{exercise}[$\star$]\label{exo:b2:linalg:3}
Sea $F = \{(x,y,z,t) \in \R^4 : x + y = z + t \text{ y } x = 2y\}$.
Da una base de $F^\circ$ y comprueba el~\cref{thm:b2:linalg:annihilator} sobre las dimensiones.
\end{exercise}

\begin{exercise}[$\star\star$]\label{exo:b2:linalg:4}
Sean $a_0, \dots, a_n$ puntos distintos de $K$ y $\varphi_i \colon P
\mapsto P(a_i)$ sobre $K_n[X]$. Demuestra que $(\varphi_0, \dots,
\varphi_n)$ es una base de $K_n[X]^*$, identifica su base
predual y desarrolla la forma $P \mapsto \int_0^1 P(t)\,\dd t$
(para $K = \R$, $n = 2$, $a_i = 0, \frac12, 1$) en esa base,
reconociendo la regla de Simpson.
\end{exercise}

\begin{exercise}[$\star\star$]\label{exo:b2:linalg:5}
Sea $u \in \mathcal{L}(E)$ con $\dim E = n$ y $\operatorname{rk} u =
1$. Demuestra que $u = \varphi(\cdot)\, a$ para cierto vector $a$ y
cierta forma $\varphi$; que $\operatorname{tr} u = \varphi(a)$; y que
$u^2 = (\operatorname{tr} u)\, u$. Deduce $\det(I + u) = 1 +
\operatorname{tr} u$.
\end{exercise}

\begin{exercise}[$\star\star$]\label{exo:b2:linalg:6}
Demuestra que todo hiperplano de $\mathcal{M}_n(K)$ ($n \geq 2$)
contiene una matriz invertible. \emph{Indicación: un hiperplano es
$\{M : \operatorname{tr}(AM) = 0\}$ para cierta $A \neq 0$
(\cref{exo:b2:linalg:9}). Si $A$ es escalar, exhibe una matriz
invertible de traza nula; en caso contrario, halla una $M$ invertible
que haga que $AM$ tenga diagonal nula ---una matriz de tipo
permutación sirve---.}
\end{exercise}

\begin{exercise}[$\star\star$]\label{exo:b2:linalg:7}
(Derivada del \omterm{def:b2:linalg:det}{determinante}) Para $A \in \mathcal{M}_n(\R)$, demuestra
a partir de la multilinealidad que
\[
\frac{\dd}{\dd t}\Big|_{t=0} \det(I_n + tA) = \operatorname{tr} A ,
\]
y deduce $\det(\eu^{tA}) = \eu^{t\operatorname{tr} A}$ suponiendo la
derivabilidad de $t \mapsto \det(\eu^{tA})$ y la propiedad de grupo
$\eu^{(s+t)A} = \eu^{sA}\eu^{tA}$ (establecida en el~\cref{ch:b2:diffeq}).
\end{exercise}

\begin{exercise}[$\star\star$]\label{exo:b2:linalg:8}
(Circulante $3 \times 3$) Sea $j = \eu^{2\iu\pi/3}$ y
\[
C = \begin{pmatrix}
a & b & c\\
c & a & b\\
b & c & a
\end{pmatrix} \in \mathcal{M}_3(\C).
\]
Comprueba que las columnas de la matriz de Vandermonde de $1, j, j^2$
son vectores propios de $C$ y deduce
\[
\det C = (a + b + c)(a + bj + cj^2)(a + bj^2 + cj).
\]
\end{exercise}

\begin{exercise}[$\star\star\star$]\label{exo:b2:linalg:9}
Demuestra que toda forma lineal $t$ sobre $\mathcal{M}_n(K)$ es $M
\mapsto \operatorname{tr}(AM)$ para una única $A$: la aplicación $A
\mapsto \operatorname{tr}(A\,\cdot)$ es un isomorfismo de
$\mathcal{M}_n(K)$ sobre su \omterm{def:b2:linalg:dual}{dual}. Deduce de nuevo el enunciado de
unicidad de la~\cref{prop:b2:linalg:trace}.
\end{exercise}

\begin{exercise}[$\star\star\star$]\label{exo:b2:linalg:10}
Sean $u, v \in \mathcal{L}(E)$ con $u \circ v - v \circ u = u$.
Demuestra que $u$ es nilpotente. \emph{Indicación: prueba que
$\operatorname{tr}(u^k) = 0$ para todo $k \geq 1$ (calcula $u^k v - v
u^k$ por inducción) y usa después el hecho siguiente, que puede
demostrarse con las identidades de Newton o por inducción sobre la
dimensión: un endomorfismo de un $\C$-espacio vectorial cuyas
potencias tienen todas traza nula es nilpotente. Trabaja sobre $\C$.}
\end{exercise}

\begin{exercise}[$\star\star$]\label{exo:b2:linalg:11}
(Vandermonde) Para $a_0, \dots, a_n \in K$, demuestra
\[
\det\begin{pmatrix}
1 & 1 & \cdots & 1\\
a_0 & a_1 & \cdots & a_n\\
\vdots & \vdots & & \vdots\\
a_0^n & a_1^n & \cdots & a_n^n
\end{pmatrix}
= \prod_{0 \leq i < j \leq n} (a_j - a_i).
\]
\emph{(Considera el \omterm{def:b2:linalg:det}{determinante} como polinomio en $a_n$: identifica
su grado, sus raíces y su coeficiente director; después induce.)}
\end{exercise}

\begin{exercise}[$\star\star\star$]\label{exo:b2:linalg:12}
Sean $A, B, C, D \in \mathcal{M}_n(K)$ con $K$ infinito, y supongamos
$CD = DC$. Demuestra que
\[
\det\begin{pmatrix} A & B\\ C & D\end{pmatrix}
= \det(AD - BC).
\]
\emph{(Trata primero el caso $D$ invertible, multiplicando por la
derecha por $\begin{psmallmatrix} I & 0\\ -D^{-1}C &
I\end{psmallmatrix}$; después sustituye $D$ por $D + tI$ y compara
dos polinomios en $t$.)}
\end{exercise}

\section{Problema: la alternativa de Fredholm}

¿Cuándo tiene solución el sistema lineal $u(x) = b$? La respuesta
completa es un enunciado de dualidad: \emph{exactamente cuando $b$ es
anulado por toda forma lineal que anula la imagen de $u$}, y esas
formas son calculables, pues son el núcleo de la \omterm{def:b2:linalg:transpose}{traspuesta}. Este
problema de fin de semana construye el diccionario completo de la
dualidad en dimensión finita (factorización de formas, bidualidad,
cálculo de \omterm{def:b2:linalg:annihilator}{anuladores}, \omterm{def:b2:linalg:transpose}{traspuesta}), demuestra la \emph{alternativa
de Fredholm} en dimensión finita y termina con la forma traza y una
caracterización: la traza es el único invariante lineal de la
semejanza. En todo el problema, $E$ y $F$ son $K$-espacios vectoriales
de dimensión finita y $n = \dim E$.

\begin{problem}[{Problema de fin de semana --- la dualidad en
dimensión finita y la alternativa de Fredholm}]
\label{pb:b2:linalg:1}
Notación: para $S \subseteq E^*$, el \emph{preanulador} es $S_\circ
= \{x \in E : \varphi(x) = 0 \text{ para toda } \varphi \in S\}$;
los \omterm{def:b2:linalg:annihilator}{anuladores} $F^\circ$ y las \omterm{def:b2:linalg:transpose}{traspuestas} $u^{\mathsf T}$ son
los de la~\cref{def:b2:linalg:annihilator} y de la~\cref{def:b2:linalg:transpose}.

\medskip
\noindent\textbf{Parte I --- El lema de factorización.} Sean
$\varphi_1, \dots, \varphi_p, \varphi \in E^*$.
\begin{enumerate}
  \item Sea $\Phi \colon E \to K^p$, $x \mapsto (\varphi_1(x),
        \dots, \varphi_p(x))$. Identifica $\ker\Phi$, prueba que
        $\Phi^{\mathsf T}$ envía las formas coordenadas de $K^p$ a
        los $\varphi_i$, y deduce
        \[
        \dim \bigl(\ker\varphi_1 \cap \dots \cap
        \ker\varphi_p\bigr) = n - \dim
        \operatorname{Vect}(\varphi_1, \dots, \varphi_p).
        \]
  \item (Lema de factorización) Demuestra la equivalencia:
        \[
        \varphi \in \operatorname{Vect}(\varphi_1, \dots,
        \varphi_p)
        \iff
        \ker\varphi_1 \cap \dots \cap \ker\varphi_p \subseteq
        \ker\varphi .
        \]
  \item Deduce que $(\varphi_1, \dots, \varphi_p)$ es libre si y
        solo si $\bigcap_i \ker\varphi_i$ tiene dimensión $n - p$;
        y que un subespacio de codimensión $p$ es intersección de
        $p$ hiperplanos, nunca de menos.
  \item En $\R^4$, sean $\varphi_1 = x + y - z$, $\varphi_2 = y + z
        - t$, $\psi = x + 2y - t$ y $\psi' = x + y + t$. Decide,
        mediante el lema de factorización, si $\psi$ y $\psi'$
        pertenecen a $\operatorname{Vect}(\varphi_1, \varphi_2)$.
  \item Sobre $E = \R_2[X]$, prueba que $\psi_0 \colon P \mapsto
        P(0)$, $\psi_1 \colon P \mapsto P(1)$ y $\psi_2 \colon P
        \mapsto \int_0^1 P(t)\dd t$ forman una base de $E^*$,
        calcula la base $(P_0, P_1, P_2)$ de $E$ de la que es la
        \omterm{def:b2:linalg:dual}{dual}, y halla el único $P \in \R_2[X]$ con $P(0) = 1$,
        $P(1) = 2$, $\int_0^1 P = \frac32$.
\end{enumerate}

\medskip
\noindent\textbf{Parte II --- Bidualidad y cálculo de
\omterm{def:b2:linalg:annihilator}{anuladores}.}
\begin{enumerate}[resume]
  \item Prueba que la \emph{aplicación de evaluación} $J \colon E
        \to E^{**}$, $J(x)(\varphi) = \varphi(x)$, es lineal e
        inyectiva, y por tanto un isomorfismo en dimensión finita.
  \item (Doble \omterm{def:b2:linalg:annihilator}{anulador}) Prueba que $J(F) = F^{\circ\circ} :=
        (F^\circ)^\circ$ para todo subespacio $F \subseteq E$:
        mediante la identificación $J$, el \omterm{def:b2:linalg:annihilator}{anulador} del \omterm{def:b2:linalg:annihilator}{anulador}
        es el propio subespacio.
  \item Demuestra el cálculo de \omterm{def:b2:linalg:annihilator}{anuladores}: $(F + G)^\circ =
        F^\circ \cap G^\circ$ y $(F \cap G)^\circ = F^\circ +
        G^\circ$.
  \item Deduce (y vuelve a demostrar directamente) que dos formas
        no nulas con el mismo núcleo son proporcionales.
  \item (Base antedual) Prueba que para toda base $(\varphi_1,
        \dots, \varphi_n)$ de $E^*$ existe una única base $(u_1,
        \dots, u_n)$ de $E$ con $\varphi_i(u_j) = \delta_{ij}$.
\end{enumerate}

\medskip
\noindent\textbf{Parte III --- El cálculo con \omterm{def:b2:linalg:transpose}{traspuestas}.}
\begin{enumerate}[resume]
  \item Prueba que $u \mapsto u^{\mathsf T}$ es una biyección
        lineal de $\mathcal{L}(E, F)$ sobre $\mathcal{L}(F^*,
        E^*)$, y que $(u^{-1})^{\mathsf T} = (u^{\mathsf T})^{-1}$
        cuando $u$ es invertible.
  \item (Naturalidad) Prueba que $u^{\mathsf T\mathsf T} \circ J_E
        = J_F \circ u$: mediante los isomorfismos de evaluación, la
        doble \omterm{def:b2:linalg:transpose}{traspuesta} \emph{es} $u$.
  \item Prueba que $u$ es sobreyectiva si y solo si $u^{\mathsf T}$
        es inyectiva, y que $u$ es inyectiva si y solo si
        $u^{\mathsf T}$ es sobreyectiva.
  \item Para $u \in \mathcal{L}(E)$: un subespacio $F$ es estable
        por $u$ si y solo si $F^\circ$ es estable por $u^{\mathsf
        T}$.
  \item Prueba que $\ker(u^{\mathsf T} - \lambda\,
        \mathrm{id}_{E^*}) = \bigl(\operatorname{im}(u - \lambda\,
        \mathrm{id}_E)\bigr)^\circ$, y deduce que $u$ y $u^{\mathsf
        T}$ tienen los mismos valores propios con las mismas
        multiplicidades geométricas.
\end{enumerate}

\medskip
\noindent\textbf{Parte IV --- La alternativa de Fredholm.}
\begin{enumerate}[resume]
  \item Demuestra que $\operatorname{im} u = (\ker u^{\mathsf
        T})_\circ$ para $u \in \mathcal{L}(E, F)$, y deduce la
        \emph{alternativa de Fredholm} en dimensión finita: la
        ecuación $u(x) = b$ tiene solución si y solo si toda $\psi
        \in F^*$ con $u^{\mathsf T}\psi = 0$ cumple $\psi(b) = 0$.
  \item Forma matricial: para $A \in \mathcal{M}_{m,n}(K)$ y $b \in
        K^m$ se cumple exactamente una de las dos afirmaciones
        siguientes: (i) $Ax = b$ tiene solución; (ii) existe $y \in
        K^m$ con $A^{\mathsf T}y = 0$ e $y^{\mathsf T}b = 1$.
        Demuestra tanto el ``a lo sumo una'' como el ``al menos
        una''.
  \item Halla todos los $b \in \R^3$ para los que el sistema
        \[
        x + y = b_1, \qquad y + z = b_2, \qquad x + 2y + z = b_3
        \]
        tiene solución, calculando el núcleo de la matriz
        \omterm{def:b2:linalg:transpose}{traspuesta}.
  \item (Un problema de Neumann discreto) Sobre $E = \R^n$ ($n \geq
        3$), definamos $L$ por $(Lx)_k = x_k - \frac12(x_{k-1} +
        x_{k+1})$, con índices módulo $n$. Prueba que $L^{\mathsf T}
        = L$ (con las identificaciones canónicas), que $\ker L$ es
        la recta de los vectores constantes \emph{(examina una
        coordenada máxima)}, y concluye: $Lx = b$ tiene solución si
        y solo si $\sum_k b_k = 0$.
\end{enumerate}

\medskip
\noindent\textbf{Parte V --- La forma traza y el teorema de
invariancia.} Recordemos del~\cref{exo:b2:linalg:9} que $A \mapsto
\operatorname{tr}(A\,\cdot)$ identifica $\mathcal{M}_n(K)$ con su
\omterm{def:b2:linalg:dual}{dual}. Supongamos $\operatorname{char} K = 0$ (por ejemplo, $K =
\Q, \R, \C$).
\begin{enumerate}[resume]
  \item Con esa identificación, prueba que el \omterm{def:b2:linalg:annihilator}{anulador} del
        subespacio $\mathcal{S}_n$ de las matrices simétricas es el
        subespacio $\mathcal{A}_n$ de las antisimétricas, y
        recíprocamente.
  \item Prueba que el \omterm{def:b2:linalg:annihilator}{anulador} del hiperplano $\mathfrak{sl}_n =
        \{M : \operatorname{tr} M = 0\}$ es la recta $K I_n$;
        equivalentemente, que toda forma lineal que se anula sobre
        las matrices de traza nula es múltiplo de la traza.
  \item Prueba que toda matriz de $\mathcal{M}_n(K)$ es suma de dos
        matrices invertibles.
  \item (La traza es el único invariante lineal de la semejanza)
        Sea $t$ una forma lineal sobre $\mathcal{M}_n(K)$ con
        $t(PMP^{-1}) = t(M)$ para toda $M$ y toda $P$ invertible.
        Prueba primero que $t(PX) = t(XP)$ para $P$ invertible,
        después que $t(BX) = t(XB)$ para \emph{toda} $B$, y
        concluye que $t = c \operatorname{tr}$ para cierto $c \in
        K$.
  \item Prueba que $\operatorname{rk} u \leq r$ si y solo si $u$ es
        suma de $r$ aplicaciones de rango $\leq 1$, es decir, $u =
        \sum_{i=1}^{r} \psi_i(\cdot)\,f_i$ con $\psi_i \in E^*$,
        $f_i \in F$; deduce $\operatorname{rk}(u + v) \leq
        \operatorname{rk} u + \operatorname{rk} v$.
  \item (Síntesis) Redacta el diccionario demostrado en este
        problema: subespacios frente a \omterm{def:b2:linalg:annihilator}{anuladores}, sumas frente a
        intersecciones, aplicaciones frente a \omterm{def:b2:linalg:transpose}{traspuestas},
        resolubilidad frente a ortogonalidad al núcleo \omterm{def:b2:linalg:transpose}{traspuesto},
        traza frente a semejanza. Cita para cada entrada la
        pregunta que la demostró y di en una frase qué sustituye a
        los recuentos de dimensiones cuando la dimensión pasa a ser
        infinita (el volumen del tercer año lo precisa sobre
        espacios de Hilbert).
\end{enumerate}
\end{problem}
